\documentclass[11pt,a4paper]{article}
\usepackage[utf8]{inputenc}
\usepackage[T1]{fontenc}
\usepackage[french]{babel}
\usepackage[a4paper,top=2cm,bottom=2.1cm,left=2.05cm,right=2.05cm]{geometry}
\usepackage{amsmath,amssymb}
\usepackage{array}
\usepackage{booktabs}
\usepackage{enumitem}
\usepackage{eurosym}
\usepackage{fancyhdr}
\usepackage{hyperref}
\usepackage{inconsolata}
\usepackage{listings}
\usepackage{microtype}
\usepackage{tabularx}
\usepackage[most]{tcolorbox}
\usepackage{titlesec}
\usepackage{xcolor}

\definecolor{ink}{HTML}{202124}
\definecolor{muted}{HTML}{5F6368}
\definecolor{line}{HTML}{DADCE0}
\definecolor{paper}{HTML}{F8F9FA}
\definecolor{accent}{HTML}{8A1538}
\definecolor{accentsoft}{HTML}{F8E8EE}
\definecolor{green}{HTML}{0B6B57}
\definecolor{greensoft}{HTML}{E6F4EF}
\definecolor{amber}{HTML}{9A5B00}
\definecolor{ambersoft}{HTML}{FFF4D6}
\definecolor{blue}{HTML}{245B8E}
\definecolor{bluesoft}{HTML}{EAF2FA}
\definecolor{violet}{HTML}{5C4D7D}
\definecolor{violetsoft}{HTML}{F1EEF8}

\hypersetup{
  colorlinks=true,
  linkcolor=accent,
  urlcolor=green
}

\setlength{\parindent}{0pt}
\setlength{\parskip}{5pt}
\setlist[itemize]{leftmargin=1.35em,itemsep=2pt,topsep=2pt}
\setlist[enumerate]{leftmargin=1.55em,itemsep=3pt,topsep=3pt}
\setlist[enumerate,1]{label=\textbf{\alph*.}}

\pagestyle{fancy}
\fancyhf{}
\fancyhead[L]{Python pour économistes}
\fancyhead[R]{Université de Lille -- L2 Économie}
\fancyfoot[C]{\thepage}
\renewcommand{\headrulewidth}{0.35pt}
\setlength{\headheight}{14.5pt}

\titleformat{\section}
  {\Large\bfseries\color{accent}}
  {}{0pt}{}
  [\vspace{-0.25em}\color{line}\titlerule]
\titleformat{\subsection}
  {\large\bfseries\color{ink}}
  {}{0pt}{}

\lstdefinelanguage{terminal}{
  morekeywords={ls,cd,pwd,mkdir,python,python3,code,dir,type,echo},
  sensitive=true
}

\lstset{
  basicstyle=\ttfamily\small,
  backgroundcolor=\color{paper},
  frame=single,
  framerule=0.35pt,
  rulecolor=\color{line},
  language=Python,
  keywordstyle=\color{accent}\bfseries,
  commentstyle=\color{muted},
  stringstyle=\color{green},
  showstringspaces=false,
  breaklines=true,
  columns=fullflexible,
  keepspaces=true,
  upquote=true,
  literate=
   {à}{{\`a}}1 {â}{{\^a}}1 {ä}{{\"a}}1
   {ç}{{\c c}}1
   {é}{{\'e}}1 {è}{{\`e}}1 {ê}{{\^e}}1 {ë}{{\"e}}1
   {î}{{\^\i}}1 {ï}{{\"\i}}1
   {ô}{{\^o}}1 {ö}{{\"o}}1
   {ù}{{\`u}}1 {û}{{\^u}}1 {ü}{{\"u}}1
   {À}{{\`A}}1 {Â}{{\^A}}1 {Ä}{{\"A}}1
   {Ç}{{\c C}}1
   {É}{{\'E}}1 {È}{{\`E}}1 {Ê}{{\^E}}1 {Ë}{{\"E}}1
   {Î}{{\^I}}1 {Ï}{{\"I}}1
   {Ô}{{\^O}}1 {Ö}{{\"O}}1
   {Ù}{{\`U}}1 {Û}{{\^U}}1 {Ü}{{\"U}}1
   {œ}{{\oe}}1 {Œ}{{\OE}}1
   {–}{{--}}1 {—}{{---}}1
}

\newcommand{\code}[1]{\texttt{#1}}
\newcommand{\pp}{\,points de pourcentage}

\newcounter{exercise}
\newcounter{extension}
\newcounter{logic}

\newcommand{\workbooktitle}[4]{%
  \setcounter{exercise}{0}%
  \setcounter{extension}{0}%
  \setcounter{logic}{0}%
  \fancyhead[L]{#1 -- #2}%
  {\Large\bfseries\color{ink}#1 -- #2\par}
  \vspace{0.2em}
  {\color{muted}Python pour économistes -- Université de Lille, L2 Économie\par}
  \vspace{0.65em}
  {\color{accent}\rule{\linewidth}{0.8pt}}
  \vspace{0.5em}
  \begin{routebox}{Question fil rouge}
  #3
  \end{routebox}
  \begin{deliverablebox}{Livrables}
  #4
  \end{deliverablebox}
  \begin{methodbox}{Méthode de travail}
  \begin{tabularx}{\linewidth}{@{}>{\bfseries\color{blue}}p{0.18\linewidth}X@{}}
  Lire & identifier les données, la question et le livrable attendu.\\
  Décomposer & écrire les étapes en pseudo-code avant de coder.\\
  Exécuter & lancer le script complet depuis le bon dossier.\\
  Vérifier & contrôler les sorties, les fichiers créés et l’interprétation.
  \end{tabularx}
  \end{methodbox}
}

\newtcolorbox{routebox}[1]{
  enhanced,breakable,
  title={#1},
  colback=paper,
  colframe=line,
  coltitle=ink,
  fonttitle=\bfseries,
  boxrule=0.5pt,
  arc=2pt,
  left=8pt,right=8pt,top=7pt,bottom=7pt,
  borderline west={2.4pt}{0pt}{accent}
}

\newtcolorbox{deliverablebox}[1]{
  enhanced,breakable,
  title={#1},
  colback=greensoft,
  colframe=green,
  coltitle=ink,
  fonttitle=\bfseries,
  boxrule=0.6pt,
  arc=2pt,
  left=8pt,right=8pt,top=7pt,bottom=7pt,
  borderline west={2.4pt}{0pt}{accent}
}


\newtcolorbox{methodbox}[1]{
  enhanced,breakable,
  title={#1},
  colback=bluesoft,
  colframe=blue,
  coltitle=ink,
  fonttitle=\bfseries,
  boxrule=0.55pt,
  arc=2pt,
  left=8pt,right=8pt,top=6pt,bottom=6pt,
  before skip=7pt,after skip=9pt,
  borderline west={2.4pt}{0pt}{blue}
}

\newtcolorbox{pathwaybox}[1]{
  enhanced,breakable,
  title={#1},
  colback=violetsoft,
  colframe=violet,
  coltitle=ink,
  fonttitle=\bfseries,
  boxrule=0.55pt,
  arc=2pt,
  left=8pt,right=8pt,top=6pt,bottom=6pt,
  before skip=8pt,after skip=8pt,
  borderline west={2.4pt}{0pt}{violet}
}

\newenvironment{pseudocodeexercise}[1]{%
  \refstepcounter{logic}%
  \begin{tcolorbox}[
    enhanced,breakable,
    title={Logique \thelogic{} -- #1 \hfill\scriptsize PSEUDO-CODE},
    colback=bluesoft,
    colframe=blue,
    coltitle=ink,
    fonttitle=\bfseries,
    boxrule=0.6pt,
    arc=2pt,
    left=8pt,right=8pt,top=6pt,bottom=6pt,
    before skip=8pt,after skip=8pt,
    borderline west={2.6pt}{0pt}{blue}
  ]%
}{\end{tcolorbox}}

\newenvironment{mandatoryexercise}[1]{%
  \refstepcounter{exercise}%
  \begin{tcolorbox}[
    enhanced,breakable,
    title={Exercice \theexercise{} -- #1 \hfill\scriptsize OBLIGATOIRE},
    colback=white,
    colframe=accent,
    coltitle=ink,
    fonttitle=\bfseries,
    boxrule=0.7pt,
    arc=2pt,
    left=8pt,right=8pt,top=7pt,bottom=7pt,
    before skip=9pt,after skip=9pt,
    borderline west={2.4pt}{0pt}{accent}
  ]%
}{\end{tcolorbox}}

\newenvironment{extensionexercise}[1]{%
  \refstepcounter{extension}%
  \begin{tcolorbox}[
    enhanced,breakable,
    title={Extension \theextension{} -- #1},
    colback=ambersoft,
    colframe=amber,
    coltitle=ink,
    fonttitle=\bfseries,
    boxrule=0.55pt,
    arc=2pt,
    left=8pt,right=8pt,top=7pt,bottom=7pt,
    before skip=8pt,after skip=8pt,
    borderline west={2.4pt}{0pt}{amber}
  ]%
}{\end{tcolorbox}}

\newtcolorbox{checkpointbox}[1]{
  enhanced,breakable,
  title={#1},
  colback=accentsoft,
  colframe=accent,
  coltitle=ink,
  fonttitle=\bfseries,
  boxrule=0.55pt,
  arc=2pt,
  left=8pt,right=8pt,top=7pt,bottom=7pt,
  borderline west={2.4pt}{0pt}{accent}
}


\begin{document}

\workbooktitle
  {TP4}
  {Graphiques et simulation simple à partir d'OWID}
  {Comment utiliser Python pour visualiser une distribution et comprendre intuitivement la stabilisation d'une moyenne ?}
  {Un fichier \code{tp4\_nom\_prenom.py}, deux figures \code{.png} et une interprétation courte.}

\begin{checkpointbox}{Parcours obligatoire}
Les exercices 1 à 8 sont obligatoires. La simulation est volontairement simple : elle sert à apprendre les boucles, \code{random.choice} et les graphiques, pas à faire de statistique avancée.
\end{checkpointbox}

\begin{checkpointbox}{Source des données}
Les taux de fécondité sont des extraits pédagogiques arrondis d'Our World in Data.
\begin{lstlisting}[language=terminal]
https://ourworldindata.org/grapher/children-per-woman-un.csv
\end{lstlisting}
\end{checkpointbox}

\begin{pseudocodeexercise}{Visualiser puis simuler}
\begin{lstlisting}
# Objectif : observer une distribution puis simuler des tirages
# Entree : liste de taux de fecondite
# 1. Calculer la moyenne observee
# 2. Tracer un graphique en barres
# 3. Tirer au hasard une valeur dans la liste
# 4. Repeter le tirage plusieurs fois
# 5. Calculer la moyenne des tirages
# 6. Observer si la moyenne se stabilise quand n augmente
\end{lstlisting}
\end{pseudocodeexercise}

\section*{A. Préparer les données et un graphique}

\begin{mandatoryexercise}{Créer l'extrait OWID}
Créer un dictionnaire où les clés sont des pays ou régions et les valeurs des taux de fécondité arrondis.
\begin{lstlisting}
fertilite = {
    "World": 2.3,
    "France": 1.8,
    "United States": 1.6,
    "Japan": 1.2,
    "India": 2.0,
    "Nigeria": 4.5,
    "Niger": 6.1,
    "Brazil": 1.6,
    "Ethiopia": 3.9,
    "China": 1.0
}
\end{lstlisting}
Afficher le nombre d'observations et la moyenne simple de ces valeurs.
\end{mandatoryexercise}

\begin{mandatoryexercise}{Graphique en barres}
Avec \code{matplotlib}, tracer un graphique en barres des taux de fécondité.
\begin{lstlisting}
import matplotlib.pyplot as plt

pays = list(fertilite.keys())
valeurs = list(fertilite.values())

plt.bar(pays, valeurs)
plt.xticks(rotation=45, ha="right")
plt.ylabel("Enfants par femme")
plt.title("Taux de fecondite - extrait OWID")
plt.tight_layout()
plt.savefig("fertilite_barres.png", dpi=300)
\end{lstlisting}
\end{mandatoryexercise}

\begin{mandatoryexercise}{Interpréter le graphique}
Ajouter trois commentaires :
\begin{enumerate}
  \item quel pays a le taux le plus élevé dans l'extrait ?
  \item quel pays a le taux le plus faible ?
  \item pourquoi ce graphique est-il plus lisible qu'une liste brute ?
\end{enumerate}
\end{mandatoryexercise}

\section*{B. Simulation très guidée}

\begin{mandatoryexercise}{Un tirage aléatoire}
Importer \code{random}, puis tirer une valeur au hasard dans la liste des taux.
\begin{lstlisting}
import random

tirage = random.choice(valeurs)
print(tirage)
\end{lstlisting}
Exécuter plusieurs fois le script : le résultat est-il toujours le même ?
\end{mandatoryexercise}

\begin{mandatoryexercise}{Moyenne de plusieurs tirages}
Compléter le code suivant pour calculer la moyenne de 20 tirages.
\begin{lstlisting}
tirages = []
for i in range(20):
    tirages.append(random.choice(valeurs))

moyenne_tirages = sum(tirages) / len(tirages)
print(moyenne_tirages)
\end{lstlisting}
Comparer cette moyenne à la moyenne de la liste initiale.
\end{mandatoryexercise}

\begin{mandatoryexercise}{Stabilisation de la moyenne}
Calculer la moyenne des tirages pour plusieurs tailles : \code{5}, \code{20}, \code{100}, \code{1000}. Afficher les résultats dans une phrase.
\begin{lstlisting}
tailles = [5, 20, 100, 1000]
for n in tailles:
    tirages = []
    for i in range(n):
        tirages.append(random.choice(valeurs))
    print(n, sum(tirages) / len(tirages))
\end{lstlisting}
\end{mandatoryexercise}

\section*{C. Visualiser la simulation}

\begin{mandatoryexercise}{Graphique de stabilisation}
Créer deux listes : \code{tailles} et \code{moyennes\_simulees}. Tracer \code{moyennes\_simulees} en fonction de \code{tailles}. Ajouter une ligne horizontale correspondant à la moyenne de la liste initiale avec \code{plt.axhline(...)}. Sauvegarder \code{simulation\_moyenne.png}.
\end{mandatoryexercise}

\begin{mandatoryexercise}{Interprétation économique}
Ajouter quatre commentaires :
\begin{enumerate}
  \item que représente la moyenne de l'extrait ?
  \item que montre la simulation quand le nombre de tirages augmente ?
  \item pourquoi ce résultat ne décrit-il pas directement la population mondiale ?
  \item pourquoi la sélection des pays est-elle importante ?
\end{enumerate}
\end{mandatoryexercise}

\section*{D. Entraînement facultatif}

\begin{extensionexercise}{Pour aller plus loin, sans obligation}
Choisir au moins deux questions si le parcours obligatoire est terminé.
\begin{enumerate}
  \item Fixer la graine avec \code{random.seed(123)} et observer l'effet.
  \item Faire 10 simulations de taille 20 et comparer les moyennes obtenues.
  \item Tracer un histogramme des 1000 tirages.
  \item Supprimer \code{Niger} de la liste et recalculer la moyenne.
  \item Ajouter une ligne pour la médiane de la liste initiale.
  \item Écrire une fonction \code{moyenne\_tirages(n)}.
  \item Tester des tailles \code{10}, \code{50}, \code{500}, \code{5000}.
  \item Calculer l'écart entre moyenne simulée et moyenne initiale.
  \item Lire intuitivement la loi des grands nombres : que veut dire “stabiliser” ?
  \item Extension avancée : construire un intervalle \code{moyenne +/- 2 * ecart\_type / racine(n)} avec une formule donnée.
\end{enumerate}
\end{extensionexercise}

\end{document}
